\Needspace{15\baselineskip}
\section*{De las secciones supervivientes a las constantes fundamentales}
\addcontentsline{toc}{chapter}{De las secciones supervivientes a las constantes fundamentales}

El continuo genealógico ya construido no publica meros desarrollos decimales: publica secciones compatibles que conservan orientación, frontera, acarreo y memoria. Una constante fundamental es la evaluación estable de una de esas secciones bajo un lector tipado. La Parte III reconstruye por ello el objeto, el mapa y la prueba que preceden a cada valor; las comparaciones arquimedianas o metrológicas aparecen sólo después de la salida interna.

En el sentido fijado por la
\hyperref[sec:tesis-inaugural-helicoidal-hmt-md]{tesis inaugural}, la lectura
nonádica no vuelve a actuar como umbral ni reinicia el continuo. La conexión nonádica conjunta y la prolongación
\[
w_6\longrightarrow w_{12}\longrightarrow w_{18}\longrightarrow w_{24}\longrightarrow w_{30}\longrightarrow R_{36}\longrightarrow G_9
\]
constituyen, con los cambios e invariantes fijados en la
\hyperref[sec:umbrales-prolongacion-tpk-inaugural]{arquitectura de umbrales
TPK}, el primer tramo certificado de la ley uniforme que construye un único
estado coinductivo enriquecido. Entre sus salidas correlacionadas, que forman
el núcleo mínimo de su imagen nonádica tipada, se encuentran
\(\pi_{\rm HMT},e_{\rm HMT},\varphi_{\rm HMT}\) y \(\alpha_{\rm HMT}\). La unidad de origen no aplana sus tipos: clausura angular, propagación, autoescala y acoplamiento poseen lectores propios; la coordenada \(\alpha\) se publica mediante el registro firmado, el sello reversible y el cierre dodecafásico. Sólo después de publicadas esas salidas se abren, con dominios separados, la realización compleja, la memoria modular, la escala de acción y la geometría angular. El estado común no queda agotado por estas cuatro coordenadas: sus lectores tipados publican además invariantes, agregados, relaciones especulares y constantes matemáticas y físicas derivadas.
